\section{Conclusion}
\label{sec:conclusion}

We presented Murmur, a system for automated prediction strategy discovery that combines domain-aware context injection, evolutionary code search, and agent-based market simulation. Our experiments on Polymarket demonstrate that context enrichment provides substantial accuracy improvements in information-sensitive domains (8--11\% in awards and commodity predictions) while confirming that high-liquidity markets such as sports remain efficiently priced and resistant to LLM-augmented analysis.

The Evolution Engine discovers interpretable strategies that match market-level Brier Scores through code-level mutation with single-change constraints and diversity pressure, avoiding the overfitting and complexity explosion that plague unconstrained approaches. Our analysis of prediction alpha across eleven domains identifies the structural conditions---information asymmetry, reasoning advantage, and low liquidity---under which systematic prediction edge can exist, providing a framework for practitioners to assess where automated forecasting tools are likely to add value.

Looking forward, three directions merit investigation. First, the Virtual Prediction Market requires comprehensive evaluation with adversarial agent training to assess strategy robustness under dynamic market conditions. Second, cross-market correlation exploitation---where information about one market informs predictions in related markets---represents an untapped source of potential alpha. Third, the architectural generality demonstrated across prediction markets, supply chain optimization, and healthcare coordination suggests that context-enriched agent-based reasoning may constitute a broadly applicable paradigm for decision-making under uncertainty, warranting systematic comparative study across domains.

\begin{acks}
The authors thank the Polymarket team for maintaining open API access to market data. Computational resources were provided by an NVIDIA DGX Spark. G.~Zhang and K.~Zhao contributed as independent researchers; the work was conducted independently and does not represent the views of any employer.
\end{acks}
